\documentclass[11pt]{amsart}
\usepackage{amsmath,amssymb,amsthm}
\usepackage[margin=1.2in]{geometry}
\usepackage{booktabs}
\usepackage{hyperref}

\newtheorem{theorem}{Theorem}[section]
\newtheorem{lemma}[theorem]{Lemma}
\newtheorem{proposition}[theorem]{Proposition}
\newtheorem{corollary}[theorem]{Corollary}
\newtheorem{conjecture}[theorem]{Conjecture}
\newtheorem{problem}[theorem]{Problem}
\theoremstyle{definition}
\newtheorem{definition}[theorem]{Definition}
\newtheorem{remark}[theorem]{Remark}
\newtheorem{observation}[theorem]{Observation}

\newcommand{\ZZ}{\mathbb{Z}}

% TODO(Tyler): author list is a placeholder; add coauthors as appropriate.

\title[Vertical supports and the hypercube tree formula]{Spanning trees of
$G\times K_2$ by vertical support, partial doubles,\\ and the bijective
program for the hypercube tree formula}
\author{Tyler Aden}
\date{Draft of August 19, 2026}

\begin{document}
\maketitle

\begin{abstract}
For a connected graph $G$, we count the spanning trees of the prism
$G\times K_2$ by their \emph{vertical support} --- the set $S$ of vertices
of $G$ whose vertical edge is used: the count is
$\tfrac12\,c(G)\,2^{|S|}F_G(S)$, where $c(G)$ is the number of spanning
trees and $F_G(S)$ the number of $S$-rooted spanning forests of $G$.
Iterating this identity over the $n$ directions of the hypercube
$Q_n = Q_{n-1}\times K_2$ telescopes, through Pascal's rule, to Stanley's
product formula $c(Q_n)=2^{2^n-n-1}\prod_k k^{\binom nk}$, giving an
elementary inductive derivation of that formula and, more importantly,
reducing the problem of a fully bijective proof --- open since
\cite[Example~5.6.10]{ec2}, and still open in bijective form after
Bernardi's combinatorial proofs \cite{bernardi} --- to a single per-level
bijective step. We prove the full-support slice of that step by an explicit
bijection, reformulate the general slice as a spanning-tree identity for
\emph{partial doubles} of $G$, report an exactly verified spin-uniformity
phenomenon (Bernardi's independence holds exactly for rooted-anywhere trees
and provably fails for root-fixed ones), and record the experimental
fine structure of the remaining step at the level of level-projections,
including where its naive product law breaks.
\end{abstract}

\section{Background}

Stanley's \emph{Enumerative Combinatorics} vol.~2 derives
\begin{equation}\label{eq:stanley}
  c(Q_n) \;=\; 2^{2^n-n-1}\prod_{k=1}^{n} k^{\binom nk}
\end{equation}
from the Laplacian eigenvalues of $Q_n$ and remarks (1st ed., p.~62,
closing Example~5.6.10) that ``a direct combinatorial proof of this formula
is not known.'' Bernardi \cite{bernardi} gave two combinatorial proofs ---
one via an independence property of edge-direction ``spins'' in uniformly
random rooted forests, one via killing involutions --- and the second
edition of \cite{ec2} (2024, Notes to Ch.~5) updates the remark: a
\emph{bijective} proof is not known. This paper organizes a route to one.

Throughout, $G$ is a connected graph on $M$ vertices (for
\eqref{eq:stanley}, $G = Q_{n-1}$ and $M = 2^{n-1}$), and
$G\times K_2$ is the prism over $G$: two copies $G\times\{0\}$,
$G\times\{1\}$ (\emph{horizontal} edges) plus a \emph{vertical} edge
$\{(u,0),(u,1)\}$ for each $u\in V(G)$. For a spanning tree $T$ of
$G\times K_2$, let $S(T)\subseteq V(G)$ be its vertical support. For
$S \subseteq V(G)$, let $F_G(S)$ denote the number of spanning forests of
$G$ with exactly $|S|$ components, each containing one vertex of $S$; by
the all-minors Matrix-Tree theorem $F_G(S) = \det L_G[\bar S]$, the minor
of the Laplacian with the rows and columns of $S$ removed.

\section{The support-product theorem}

\begin{theorem}\label{thm:support}
For every nonempty $S\subseteq V(G)$, the number of spanning trees $T$ of
$G\times K_2$ with $S(T)=S$ is
\[
  N_G(S) \;=\; \tfrac12\, c(G)\; 2^{|S|}\, F_G(S).
\]
\end{theorem}

\begin{proof}
Weight the vertical edge at $u$ by an indeterminate $z_u$ and every
horizontal edge by $1$, and let $t(z)=\sum_T\prod_{u\in S(T)}z_u$ be the
weighted spanning-tree sum. With vertices ordered level by level, the
weighted Laplacian is
$\left(\begin{smallmatrix} L_G+Z & -Z\\ -Z & L_G+Z \end{smallmatrix}\right)$,
$Z=\operatorname{diag}(z_u)$, which the orthogonal level-swap basis
$\tfrac1{\sqrt2}\left(\begin{smallmatrix} I & I\\ I & -I
\end{smallmatrix}\right)$ block-diagonalizes as $L_G\oplus(L_G+2Z)$. For a
weighted graph on $N$ vertices the product of the nonzero Laplacian
eigenvalues is $N$ times the tree sum, and the kernel lies in the $L_G$
block; with $N = 2M$,
\[
  t(z) \;=\; \frac{1}{2M}\,\Bigl(\;\prod_{\lambda\ne0}\lambda(L_G)\Bigr)
  \det(L_G+2Z)
  \;=\; \frac{c(G)}{2}\,\det(L_G+2Z).
\]
Expanding by principal minors in the diagonal matrix $2Z$ and applying the
all-minors Matrix-Tree theorem,
$\det(L_G+2Z)=\sum_{R}2^{|R|}F_G(R)\prod_{u\in R}z_u$. Both sides are
multilinear in the $z_u$; compare coefficients of $\prod_{u\in S}z_u$.
\end{proof}

\begin{remark}
The identity is elementary and product-graph Matrix-Tree computations of
this kind appear in Martin--Reiner \cite{mr} and Bernardi
\cite[\S4]{bernardi}; we make no novelty claim for
Theorem~\ref{thm:support} itself. Its role here is architectural: it is the
exact statement whose bijective realization would prove
\eqref{eq:stanley} bijectively (Corollary~\ref{cor:telescope} and
Problem~\ref{prob:main}). We first found it as an experimental identity:
it was confirmed class-by-class for $Q_2,Q_3,Q_4$ ($3+15+255$ supports; the
$Q_4$ case over all $679\,477\,248$ tree--root pairs) before the proof was
written down.
\end{remark}

\begin{corollary}\label{cor:telescope}
Iterating Theorem~\ref{thm:support} yields \eqref{eq:stanley}.
\end{corollary}

\begin{proof}
Summing the theorem over all $S$ (the $S=\varnothing$ term vanishes) and
using $\sum_R z^{|R|}F_G(R) = \det(zI+L_G)$ at $z=2$:
\[
  c(G\times K_2) \;=\; \tfrac12\,c(G)\,\det(2I+L_G).
\]
For $G=Q_{n-1}$, whose Laplacian eigenvalues are $2k$ with multiplicity
$\binom{n-1}k$, $\det(2I+L)=\prod_{k=0}^{n-1}(2k+2)^{\binom{n-1}k}$.
Writing $R_n=2^n c(Q_n)$, the recursion reads
$R_n = R_{n-1}\prod_{j=1}^n(2j)^{\binom{n-1}{j-1}}$, and Pascal's rule
gives $R_n=\prod_{j=1}^n(2j)^{\binom nj}$ by induction, which is
\eqref{eq:stanley}.
\end{proof}

Of course, evaluating $\det(2I+L_{Q_{n-1}})$ spectrally re-uses linear
algebra; combinatorially, $\det(2I+L_G) = \sum_S 2^{|S|}F_G(S)$ counts
\emph{pairs} (an $S$-rooted spanning forest, a $\{0,1\}$-coloring of its
roots), so the recursion has the set-level form
\begin{equation}\label{eq:setform}
\begin{gathered}
  \{\text{spanning trees of } G\times K_2\}\ \times\ \{0,1\}\\
  \longleftrightarrow\quad
  \{\text{spanning trees of }G\}\ \times\
  \{\text{root-$2$-colored rooted spanning forests of } G\},
\end{gathered}
\end{equation}
and a bijective proof of \eqref{eq:stanley} needs exactly a constructive
form of \eqref{eq:setform}, level by level.

\section{Bijective fragments}

\subsection{The full-support slice}

\begin{proposition}\label{prop:full}
A spanning tree of $G\times K_2$ containing every vertical edge is the same
data as a spanning tree of $G$ together with a level in $\{0,1\}$ for each
of its $M-1$ edges: contract the vertical edges to project, and lift each
base tree edge to its recorded level to invert. Hence
$N_G(V(G)) = c(G)\,2^{M-1}$, the $S=V(G)$ case of
Theorem~\ref{thm:support}.
\end{proposition}

\subsection{Partial doubles}

\begin{definition}
For $S\subseteq V(G)$, the \emph{partial double} $G^{(S)}$ is the graph
obtained from $G$ by splitting each $u\notin S$ into $u_0,u_1$: an edge
$(u,v)$ of $G$ becomes two parallel edges when $u,v\in S$; the two edges
$(u,v_0),(u,v_1)$ when $u\in S$, $v\notin S$; and $(u_0,v_0),(u_1,v_1)$
when $u,v\notin S$.
\end{definition}

Contracting the vertical edges of a tree with support $S$ is a bijection
onto the spanning trees of $G^{(S)}$, so Theorem~\ref{thm:support} is
equivalent to
\begin{equation}\label{eq:partial}
  c\bigl(G^{(S)}\bigr) \;=\; \tfrac12\,c(G)\,2^{|S|}\,F_G(S).
\end{equation}
For $S=V(G)$ this is Proposition~\ref{prop:full} (all edges doubled); for
$S=\varnothing$ both sides vanish on connectivity grounds. Splitting one
vertex at a time interpolates between $G^{(V)}$ and $G^{(S)}$ and is a
natural inductive route to a bijection, though the ratio
$F_G(S)/F_G(S\cup u)$ is not local, so the induction must carry more
structure than a single count.

\subsection{What the spins do and do not give}

Bernardi's independence theorem \cite[Thm.~1]{bernardi} concerns uniformly
random rooted forests conditioned on their projection; transported to
spanning trees it says the \emph{spins} (for each $u\in S$, the level of
the endpoint of the vertical edge at $u$ nearer the root) should be free
bits. We verified computationally: for the ensemble of (tree, root) pairs
of $Q_3$, conditioned on the direction and the support $S$, the joint
distribution of the spins is \emph{exactly} uniform on $\{0,1\}^S$ (all
$45$ conditioned classes, zero discrepancy), and remains consistent with
uniform on samples of $2\times10^6$ ($Q_4$) and $10^6$ ($Q_5$) uniform
trees. For trees with a \emph{fixed} root the uniformity provably fails
(e.g.\ $\chi^2=330.7$ on $65$ dof already for $Q_3$). Thus the factor
$2^{|S|}$ of Theorem~\ref{thm:support} is realized by free bits only in the
rooted-anywhere ensemble; any bijection must fix the root only after the
bits are extracted. This matches the shape of \eqref{eq:setform}, which is
stated for a $2$-fold cover of the tree set rather than for rooted trees.

\section{Fine structure of the remaining step}\label{sec:fine}

Deleting (rather than contracting) the vertical edges of a tree $T$ with
support $S$ leaves horizontal forests $H_0, H_1$ in the two levels;
projecting both to $G$ gives an edge multiset $U$ with multiplicities in
$\{1,2\}$, where a multiplicity-$1$ edge remembers a free level choice and
a multiplicity-$2$ edge uses both levels. Vertex and edge counts show that
$T$ is a tree if and only if the chosen level-lift of $(S,U)$ is connected.
Exhaustively on $Q_3$ (base $Q_2$):

\begin{itemize}
\item every occurring $U$ decomposes as (spanning tree of the base)
$\uplus$ ($S$-rooted spanning forest of the base) --- the edge budgets
match Theorem~\ref{thm:support} exactly --- and conversely every
decomposable $(S,U)$ occurs;
\item with $m$ the number of multiplicity-$1$ edges and $d$ the number of
decompositions, every class satisfies the \emph{mass formula}
$N(S,U)\cdot d = 2^m$ (the classes split as $d=1$ with all $2^m$ lifts
connected, and $d=2$ with exactly half).
\end{itemize}

On $Q_4$ (base $Q_3$; all $42\,467\,328$ trees, $2\,459\,160$ classes) the
mass formula holds in $95.7\%$ of classes; but the full data reveals a
\emph{universal law} subsuming it. In every class of $Q_3$ and of $Q_4$
without exception,
\begin{equation}\label{eq:universal}
  N(S,U) \;=\; d(S,U)\cdot 2^{\,m-2s}
  \qquad\text{for some integer } s = s(S,U) \ge 0
\end{equation}
(observed $s \le 5$ on $Q_4$): the connected level-lifts distribute over
the $d$ decompositions of $U$ in equal shares whose codimension in
$\mathbb{F}_2^{\,m}$ is always \emph{even}. The mass formula is the special
case $d = 2^s$. Probing the geometry: in a sample of $40$ exceptional
classes, the set $V \subseteq \mathbb{F}_2^m$ of connected lifts was, in
$22$ cases, exactly a union of $d$ cosets of its stabilizer subspace
$W=\{x : V\oplus x = V\}$ with $\dim W = m-2s$; in the remaining $18$, $V$
is a union of fewer, larger cosets with the decompositions clustering in
groups of $4$ per coset. We conjecture \eqref{eq:universal} for all
connected base graphs, and identifying the invariant $s(S,U)$ together with
a canonical assignment of lifts to decompositions is, in our view, the
combinatorial heart of \eqref{eq:setform}.

\section{Open problems}

\begin{problem}\label{prob:main}
Construct an explicit bijection realizing \eqref{eq:setform} ---
equivalently \eqref{eq:partial} for all $S$ --- for $G = Q_{n-1}$, or for
arbitrary connected $G$. By Corollary~\ref{cor:telescope} this yields a
bijective proof of \eqref{eq:stanley}, resolving the bijective form of the
question of \cite[Example~5.6.10]{ec2}.
\end{problem}

\begin{problem}
Prove the universal law \eqref{eq:universal} and identify the invariant
$s(S,U)$; determine when the connected lifts form exactly $d$ cosets of a
subspace of dimension $m-2s$, and construct the canonical assignment of
lifts to decompositions in general.
\end{problem}

\begin{problem}
Make Bernardi's spin-flip independence constructive on rooted-anywhere
trees (his proof is an induction through vertex deletions and arc
mergings), and determine its interaction with the contraction map of
Proposition~\ref{prop:full}.
\end{problem}

\begin{thebibliography}{9}
\bibitem{bernardi} O.~Bernardi, \emph{On the spanning trees of the
hypercube and other products of graphs}, Electron.\ J.\ Combin.\ 19(4)
(2012), \#P51.
\bibitem{ec2} R.~P.~Stanley, \emph{Enumerative Combinatorics, Volume 2},
Cambridge Univ.\ Press, 1999; 2nd ed.\ 2024.
\bibitem{mr} J.~L.~Martin, V.~Reiner, \emph{Factorization of some weighted
spanning tree enumerators}, J.\ Combin.\ Theory Ser.~A 104 (2003) 287--300.
\end{thebibliography}

\end{document}
